\section{Sharp reflection stability}\label{sec:upper}

Throughout this section let $X'$ be an independent copy of $X$, set
\[
Y=X-X',\qquad G\sim\Normal(0,2),
\]
and define
\[
F(z):=M_Y(z)-e^{z^2},\qquad
L:=\log(1/\Delta).
\]
We prove \cref{thm:sharp-upper} through five lemmas.

\subsection{High-order reflected moment transfer}

\begin{lemma}[Local defect to high-order moments]\label{lem:moment-transfer}
Fix $\rho,\tau,K$. For every $A>0$ there is $c_0=c_0(A,\rho,\tau,K)>0$ such that, with
\[
m=\left\lfloor c_0\frac{L}{\log L}\right\rfloor,
\]
all sufficiently small $\Delta$ satisfy
\[
\left|\E Y^k-\E G^k\right|
\le e^{-A m\log m},
\qquad 0\le k\le2m.
\]
In particular,
\[
\E|Y|^{2m}\le (C m)^m.
\]
\end{lemma}

\begin{proof}
Choose a fixed $0<\rho_0<\min\{\rho,\tau\}$ and a symmetric ellipse whose upper and lower half-ellipses lie in the strip $|\Re z|<2\tau$. Exponential integrability gives a bound $|F(z)|\le B_0$ on this ellipse, with $B_0$ depending only on $(\rho,\tau,K)$. On the real diameter, the defect hypothesis gives
\[
|F(s)|=e^{s^2}|e^{R_X(s)}-1|\le C_0\Delta,
\qquad |s|\le\rho_0,
\]
where $F(s)=M_Y(s)-e^{s^2}$ and $Y=X-X'$. The two-constants theorem on the upper and lower half-ellipses yields constants $r_0>0$, $\omega_0>0$, and $C_1<\infty$ such that
\[
\sup_{|z|\le r_0}|F(z)|\le C_1\Delta^{\omega_0}.
\]
The Cauchy estimate consequently gives
\[
|F^{(k)}(0)|\le C_1 k! r_0^{-k}\Delta^{\omega_0}.
\]
For $k\le2m$ and $m=c_0L/\log L$, Stirling's bound shows that the logarithm of the right-hand side is at most $-\omega_0L+2c_0L+o(L)$. Choosing $c_0$ sufficiently small in terms of the prescribed $A$ gives $|F^{(k)}(0)|\le e^{-Am\log m}$. Since $F^{(k)}(0)=\E Y^k-\E G^k$, the moment conclusion follows. Finally,
\[
\E G^{2m}=\frac{(2m)!}{m!}\le(Cm)^m,
\]
and the negligible moment mismatch gives the stated $2m$-th moment bound.
\end{proof}

\subsection{Tail lifting and bounded truncation}

\begin{lemma}[Square-root tail lifting]\label{lem:tail-lift}
There are constants $C,c>0$ with the following property. For every sufficiently large fixed $B$, if $N=B\sqrt m$, then
\[
\Pp(|X|>N)\le C e^{-\beta_Bm},
\]
where $\beta_B\to\infty$ as $B\to\infty$.
If $X_N$ denotes the conditional law $X\mid\{|X|\le N\}$, then
\[
d_{\rm TV}(X,X_N)\le C e^{-\beta_Bm}.
\]
Moreover
\[
|\E X_N|+|\operatorname{Var}(X_N)-1|
\le C(1+N^2)e^{-cN}.
\]
\end{lemma}

\begin{proof}
Let $a$ be a median of $X$. Since $\E X^2=1$, $|a|\le\sqrt2$. The safe symmetrization estimate is
\[
\Pp(|X-a|>x)\le4\Pp(|X-X'|>x).
\]
Combine this with Markov's inequality and \cref{lem:moment-transfer}. The conditional-law total-variation estimate is immediate. The mean and variance shifts are controlled separately from the original exponential moment, which gives exponentially small first and second tail moments at radius $N$.
\end{proof}

\subsection{A growing real Gaussian-product window}

\begin{lemma}[Real-axis product approximation]\label{lem:real-window}
For $B$ fixed sufficiently large, there is $a_0=a_0(B,\rho,\tau,K)>0$ such that, with
\[
R=a_0\sqrt m
\]
and $f_N$ the characteristic function of $X_N$,
\[
\sup_{|t|\le2R}
\left|f_N(t)f_N(-t)-e^{-t^2}\right|
\le e^{-\gamma_1m}
\]
for some $\gamma_1>0$.
\end{lemma}

\begin{proof}
Use the $2m-1$ Taylor polynomial of $\phi_Y(t)$ and the real-variable remainder
\[
\left|e^{iu}-\sum_{k=0}^{2m-1}\frac{(iu)^k}{k!}\right|
\le\frac{|u|^{2m}}{(2m)!}.
\]
The high-order moment mismatch from \cref{lem:moment-transfer} is negligible on $|t|\le2R$, while
\[
\frac{|t|^{2m}\E|Y|^{2m}}{(2m)!}
\le (C a_0^2)^m.
\]
Choose $a_0$ small. Finally transfer from $X$ to $X_N$ using \cref{lem:tail-lift}.
\end{proof}

\subsection{Complex propagation and a zero-free factor disk}

\begin{lemma}[Growing zero-free disk]\label{lem:zero-free}
The constants $B$ and $a_0$ can be chosen so that, for $R=a_0\sqrt m$,
\[
f_N(z)\ne0,\qquad |z|\le R.
\]
More precisely, for some $\gamma_2>a_0^2$,
\[
\sup_{|z|\le R}
\left|f_N(z)f_N(-z)-e^{-z^2}\right|
\le e^{-\gamma_2m}.
\]
\end{lemma}

\begin{proof}
Since $|X_N|\le N$, $f_N$ is entire and
\[
|f_N(z)|\le e^{N|\Im z|}.
\]
Hence, for
\[
\mathcal F_N(z):=f_N(z)f_N(-z)-e^{-z^2},
\]
one has
\[
\log|\mathcal F_N(z)|\le C(NR+R^2)\le C(Ba_0+a_0^2)m,
\qquad |z|\le2R.
\]
On the real diameter, \cref{lem:real-window} gives $\log|\mathcal F_N|\le-\gamma_1m$. Apply the two-constants theorem separately in the upper and lower half-disks of radius $2R$. The harmonic measure of the diameter is uniformly bounded below on the concentric half-disk of radius $R$. Choose first $B$ large and then $a_0$ sufficiently small so that the propagated exponent $\gamma_2$ is positive and satisfies $\gamma_2>a_0^2$. Since
\[
\inf_{|z|\le R}|e^{-z^2}|\ge e^{-R^2}=e^{-a_0^2m},
\]
relative closeness implies $f_N(z)f_N(-z)\ne0$, hence $f_N(z)\ne0$ on the disk.
\end{proof}

\subsection{Logarithmic Gaussianization and Esseen smoothing}

\begin{lemma}[Zero-free factor to Kolmogorov distance]\label{lem:factor-gaussian}
Under the conclusions of \cref{lem:zero-free},
\[
\dK(X_N,\Normal(\mu_N,\sigma_N^2))\le \frac{C}{R},
\]
where $\mu_N=\E X_N$ and $\sigma_N^2=\operatorname{Var}(X_N)$.
Consequently
\[
\dK(X,\Normal(0,1))\le \frac{C}{\sqrt m}.
\]
\end{lemma}

\begin{proof}
On $|z|<R$ choose the analytic logarithm
\[
\psi_N(z)=\log f_N(z),\qquad \psi_N(0)=0.
\]
The bounded support gives
\[
\Re\psi_N(z)=\log|f_N(z)|\le NR.
\]
Carath\'eodory's coefficient inequality therefore yields, writing $\psi_N(z)=\sum_{k\ge1}a_kz^k$,
\[
|a_k|\le2NR^{1-k}.
\]
Since $a_1=i\mu_N$ and $a_2=-\sigma_N^2/2$, for every fixed $0<\theta<1$,
\[
\psi_N(t)
=i\mu_Nt-\frac{\sigma_N^2t^2}{2}
+O\left(\frac{|t|^3}{R}\right),
\qquad |t|\le\theta R,
\]
because $N/R=B/a_0$ is fixed. By \cref{lem:tail-lift}, $\mu_N=o(R^{-1})$ and $\sigma_N^2=1+o(R^{-1})$. Taking $\theta$ sufficiently small gives Gaussian damping, and hence
\[
\left|f_N(t)-e^{i\mu_Nt-\sigma_N^2t^2/2}\right|
\le C\frac{|t|^3}{R}e^{-ct^2}.
\]
Esseen smoothing with cutoff $T=\theta R$ gives $C/R$. Add the truncation and parameter-shift errors.
\end{proof}

\begin{proof}[Proof of \cref{thm:sharp-upper}]
Apply \cref{lem:moment-transfer,lem:tail-lift,lem:real-window,lem:zero-free,lem:factor-gaussian}. Since
\[
m\asymp\frac{L}{\log L},\qquad L=\log(1/\Delta),
\]
we have
\[
\frac1{\sqrt m}
\asymp
\sqrt{\frac{\log L}{L}},
\]
which is the claimed rate.
\end{proof}
